II) \emph{Caso general}. --- Como $A'$ es un $A$-módulo sin torsión,
$A'\otimes_AK$ se identifica con un subanillo de $K'$, por lo que
$X''=\operatorname{Spec}(A'\otimes_AK)$ es un $K$-preesquema íntegro cuyo
cuerpo de funciones racionales es $K'$. Se sabe (4.6.6) que existe una
extensión radicial finita $K_1$ de $K$ tal que, si
$X''_1=\operatorname{Spec}(A'\otimes_AK_1)=X''\otimes_KK_1$,
$(X''_1)_{\mathrm{red}}$ es un $K_1$-preesquema \emph{geométricamente reducido} de
tipo finito; además, como el morfismo $\operatorname{Spec}(K_1)\to
\operatorname{Spec}(K)$ es radicial, finito y sobreyectivo, es un
homeomorfismo universal (2.4.5), de modo que $X''_1$ es homeomorfo a $X''$
y, por tanto, $(X''_1)_{\mathrm{red}}$ es íntegro; además, su cuerpo de
funciones racionales $K'_1$ es una extensión radicial finita de $K'$ y
una extensión \emph{separable} de tipo finito de $K_1$ (4.6.1). En virtud de la
hipótesis \emph{c)} del enunciado, existe una sub-$A$-álgebra \emph{finita} $A_1$ de
$K_1$, cuyo cuerpo de fracciones es $K_1$ y tal que, poniendo
$X_1=\operatorname{Spec}(A_1)$, $\operatorname{Reg}(X_1)$ contiene un abierto
no vacío de $X_1$. Sea $A'_1$ la imagen del homomorfismo canónico
$A'\otimes_AA_1\to K'_1$ y sea $X'_1=\operatorname{Spec}(A'_1)$; $A'_1$
es un anillo íntegro que es una $A'$-álgebra \emph{finita} y cuyo cuerpo de
fracciones es $K'_1$ por construcción; además, como el homomorfismo
compuesto $A'\to A'\otimes_AA_1\to A'_1\to K'_1$ es idéntico a
$A'\to K'\to K'_1$, y por tanto inyectivo, el homomorfismo
$A'\to A'_1$ es inyectivo; el morfismo $g:X'_1\to X'$ es, pues, finito y
sobreyectivo (\textbf{II}, 6.1.10). Ahora bien, la hipótesis sobre $A_1$ y la
parte I) de la demostración implican que $\operatorname{Reg}(X'_1)$
\oldpage[IV-2]{166}
contiene un abierto no vacío $V'_1$; como $g$ es cerrado (\textbf{II}, 6.1.10),
se puede suponer que $V'_1=g^{-1}(V')$, donde $V'$ es un abierto afín de
$X'$; reemplazando $A'$ por el anillo de $V'$ y $A'_1$ por el de $V'_1$,
se puede suponer, por tanto, que $X'_1$ es \emph{regular}. Además, el mismo
razonamiento que en I), aplicado al punto genérico $\xi'$ de $X'$, permite
(reemplazando, si es necesario, $X'$ por una vecindad de $\xi'$) suponer que
$A'_1$ es un $A'$-módulo \emph{libre} y, por tanto, que el morfismo $g$ es \emph{plano}.
Pero entonces se deduce de (6.5.2, (i)) que $X'$ es \emph{regular}. Queda
demostrado.
\end{proof}
